% Template for PLoS
% Version 3.5 March 2018
%
% % % % % % % % % % % % % % % % % % % % % %
%
% -- IMPORTANT NOTE
%
% This template contains comments intended
% to minimize problems and delays during our production
% process. Please follow the template instructions
% whenever possible.
%
% % % % % % % % % % % % % % % % % % % % % % %
%
% Once your paper is accepted for publication,
% PLEASE REMOVE ALL TRACKED CHANGES in this file
% and leave only the final text of your manuscript.
% PLOS recommends the use of latexdiff to track changes during review, as this will help to maintain a clean tex file.
% Visit https://www.ctan.org/pkg/latexdiff?lang=en for info or contact us at latex@plos.org.
%
%
% There are no restrictions on package use within the LaTeX files except that
% no packages listed in the template may be deleted.
%
% Please do not include colors or graphics in the text.
%
% The manuscript LaTeX source should be contained within a single file (do not use \input, \externaldocument, or similar commands).
%
% % % % % % % % % % % % % % % % % % % % % % %
%
% -- FIGURES AND TABLES
%
% Please include tables/figure captions directly after the paragraph where they are first cited in the text.
%
% DO NOT INCLUDE GRAPHICS IN YOUR MANUSCRIPT
% - Figures should be uploaded separately from your manuscript file.
% - Figures generated using LaTeX should be extracted and removed from the PDF before submission.
% - Figures containing multiple panels/subfigures must be combined into one image file before submission.
% For figure citations, please use "Fig" instead of "Figure".
% See http://journals.plos.org/plosone/s/figures for PLOS figure guidelines.
%
% Tables should be cell-based and may not contain:
% - spacing/line breaks within cells to alter layout or alignment
% - do not nest tabular environments (no tabular environments within tabular environments)
% - no graphics or colored text (cell background color/shading OK)
% See http://journals.plos.org/plosone/s/tables for table guidelines.
%
% For tables that exceed the width of the text column, use the adjustwidth environment as illustrated in the example table in text below.
%
% % % % % % % % % % % % % % % % % % % % % % % %
%
% -- EQUATIONS, MATH SYMBOLS, SUBSCRIPTS, AND SUPERSCRIPTS
%
% IMPORTANT
% Below are a few tips to help format your equations and other special characters according to our specifications. For more tips to help reduce the possibility of formatting errors during conversion, please see our LaTeX guidelines at http://journals.plos.org/plosone/s/latex
%
% For inline equations, please be sure to include all portions of an equation in the math environment.
%
% Do not include text that is not math in the math environment.
%
% Please add line breaks to long display equations when possible in order to fit size of the column.
%
% For inline equations, please do not include punctuation (commas, etc) within the math environment unless this is part of the equation.
%
% When adding superscript or subscripts outside of brackets/braces, please group using {}.
%
% Do not use \cal for caligraphic font.  Instead, use \mathcal{}
%
% % % % % % % % % % % % % % % % % % % % % % % %
%
% Please contact latex@plos.org with any questions.
%
% % % % % % % % % % % % % % % % % % % % % % % %

\documentclass[10pt,letterpaper]{article}
\usepackage[top=0.85in,left=2.75in,footskip=0.75in]{geometry}

% amsmath and amssymb packages, useful for mathematical formulas and symbols
\usepackage{amsmath,amssymb}

% Use adjustwidth environment to exceed column width (see example table in text)
\usepackage{changepage}

% Use Unicode characters when possible
\usepackage[utf8x]{inputenc}

% textcomp package and marvosym package for additional characters
\usepackage{textcomp,marvosym}

% cite package, to clean up citations in the main text. Do not remove.
% \usepackage{cite}

% Use nameref to cite supporting information files (see Supporting Information section for more info)
\usepackage{nameref,hyperref}

% line numbers
\usepackage[right]{lineno}

% ligatures disabled
\usepackage{microtype}
\DisableLigatures[f]{encoding = *, family = * }

% color can be used to apply background shading to table cells only
\usepackage[table]{xcolor}

% array package and thick rules for tables
\usepackage{array}

% create "+" rule type for thick vertical lines
\newcolumntype{+}{!{\vrule width 2pt}}

% create \thickcline for thick horizontal lines of variable length
\newlength\savedwidth
\newcommand\thickcline[1]{%
  \noalign{\global\savedwidth\arrayrulewidth\global\arrayrulewidth 2pt}%
  \cline{#1}%
  \noalign{\vskip\arrayrulewidth}%
  \noalign{\global\arrayrulewidth\savedwidth}%
}

% \thickhline command for thick horizontal lines that span the table
\newcommand\thickhline{\noalign{\global\savedwidth\arrayrulewidth\global\arrayrulewidth 2pt}%
\hline
\noalign{\global\arrayrulewidth\savedwidth}}


% Remove comment for double spacing
%\usepackage{setspace}
%\doublespacing

% Text layout
\raggedright
\setlength{\parindent}{0.5cm}
\textwidth 5.25in
\textheight 8.75in

% Bold the 'Figure #' in the caption and separate it from the title/caption with a period
% Captions will be left justified
\usepackage[aboveskip=1pt,labelfont=bf,labelsep=period,justification=raggedright,singlelinecheck=off]{caption}
\renewcommand{\figurename}{Fig}

% Use the PLoS provided BiBTeX style
% \bibliographystyle{plos2015}

% Remove brackets from numbering in List of References
\makeatletter
\renewcommand{\@biblabel}[1]{\quad#1.}
\makeatother



% Header and Footer with logo
\usepackage{lastpage,fancyhdr,graphicx}
\usepackage{epstopdf}
%\pagestyle{myheadings}
\pagestyle{fancy}
\fancyhf{}
%\setlength{\headheight}{27.023pt}
%\lhead{\includegraphics[width=2.0in]{PLOS-submission.eps}}
\rfoot{\thepage/\pageref{LastPage}}
\renewcommand{\headrulewidth}{0pt}
\renewcommand{\footrule}{\hrule height 2pt \vspace{2mm}}
\fancyheadoffset[L]{2.25in}
\fancyfootoffset[L]{2.25in}
\lfoot{\today}

%% Include all macros below

\newcommand{\lorem}{{\bf LOREM}}
\newcommand{\ipsum}{{\bf IPSUM}}

\usepackage{color}
\usepackage{fancyvrb}
\newcommand{\VerbBar}{|}
\newcommand{\VERB}{\Verb[commandchars=\\\{\}]}
\DefineVerbatimEnvironment{Highlighting}{Verbatim}{commandchars=\\\{\}}
% Add ',fontsize=\small' for more characters per line
\usepackage{framed}
\definecolor{shadecolor}{RGB}{248,248,248}
\newenvironment{Shaded}{\begin{snugshade}}{\end{snugshade}}
\newcommand{\AlertTok}[1]{\textcolor[rgb]{0.94,0.16,0.16}{#1}}
\newcommand{\AnnotationTok}[1]{\textcolor[rgb]{0.56,0.35,0.01}{\textbf{\textit{#1}}}}
\newcommand{\AttributeTok}[1]{\textcolor[rgb]{0.77,0.63,0.00}{#1}}
\newcommand{\BaseNTok}[1]{\textcolor[rgb]{0.00,0.00,0.81}{#1}}
\newcommand{\BuiltInTok}[1]{#1}
\newcommand{\CharTok}[1]{\textcolor[rgb]{0.31,0.60,0.02}{#1}}
\newcommand{\CommentTok}[1]{\textcolor[rgb]{0.56,0.35,0.01}{\textit{#1}}}
\newcommand{\CommentVarTok}[1]{\textcolor[rgb]{0.56,0.35,0.01}{\textbf{\textit{#1}}}}
\newcommand{\ConstantTok}[1]{\textcolor[rgb]{0.00,0.00,0.00}{#1}}
\newcommand{\ControlFlowTok}[1]{\textcolor[rgb]{0.13,0.29,0.53}{\textbf{#1}}}
\newcommand{\DataTypeTok}[1]{\textcolor[rgb]{0.13,0.29,0.53}{#1}}
\newcommand{\DecValTok}[1]{\textcolor[rgb]{0.00,0.00,0.81}{#1}}
\newcommand{\DocumentationTok}[1]{\textcolor[rgb]{0.56,0.35,0.01}{\textbf{\textit{#1}}}}
\newcommand{\ErrorTok}[1]{\textcolor[rgb]{0.64,0.00,0.00}{\textbf{#1}}}
\newcommand{\ExtensionTok}[1]{#1}
\newcommand{\FloatTok}[1]{\textcolor[rgb]{0.00,0.00,0.81}{#1}}
\newcommand{\FunctionTok}[1]{\textcolor[rgb]{0.00,0.00,0.00}{#1}}
\newcommand{\ImportTok}[1]{#1}
\newcommand{\InformationTok}[1]{\textcolor[rgb]{0.56,0.35,0.01}{\textbf{\textit{#1}}}}
\newcommand{\KeywordTok}[1]{\textcolor[rgb]{0.13,0.29,0.53}{\textbf{#1}}}
\newcommand{\NormalTok}[1]{#1}
\newcommand{\OperatorTok}[1]{\textcolor[rgb]{0.81,0.36,0.00}{\textbf{#1}}}
\newcommand{\OtherTok}[1]{\textcolor[rgb]{0.56,0.35,0.01}{#1}}
\newcommand{\PreprocessorTok}[1]{\textcolor[rgb]{0.56,0.35,0.01}{\textit{#1}}}
\newcommand{\RegionMarkerTok}[1]{#1}
\newcommand{\SpecialCharTok}[1]{\textcolor[rgb]{0.00,0.00,0.00}{#1}}
\newcommand{\SpecialStringTok}[1]{\textcolor[rgb]{0.31,0.60,0.02}{#1}}
\newcommand{\StringTok}[1]{\textcolor[rgb]{0.31,0.60,0.02}{#1}}
\newcommand{\VariableTok}[1]{\textcolor[rgb]{0.00,0.00,0.00}{#1}}
\newcommand{\VerbatimStringTok}[1]{\textcolor[rgb]{0.31,0.60,0.02}{#1}}
\newcommand{\WarningTok}[1]{\textcolor[rgb]{0.56,0.35,0.01}{\textbf{\textit{#1}}}}




\usepackage{forarray}
\usepackage{xstring}
\newcommand{\getIndex}[2]{
  \ForEach{,}{\IfEq{#1}{\thislevelitem}{\number\thislevelcount\ExitForEach}{}}{#2}
}

\setcounter{secnumdepth}{0}

\newcommand{\getAff}[1]{
  \getIndex{#1}{}
}

\providecommand{\tightlist}{%
  \setlength{\itemsep}{0pt}\setlength{\parskip}{0pt}}

\begin{document}
\vspace*{0.2in}

% Title must be 250 characters or less.
\begin{flushleft}
{\Large
\textbf\newline{Ten Simple Rules for Writing Dockerfiles for Reproducible Research} % Please use "sentence case" for title and headings (capitalize only the first word in a title (or heading), the first word in a subtitle (or subheading), and any proper nouns).
}
\newline
% Insert author names, affiliations and corresponding author email (do not include titles, positions, or degrees).
\\
Daniel Nüst\textsuperscript{\getAff{Institute for Geoinformatics, University of Muenster, Muenster, Germany}}\textsuperscript{*},
Vanessa Sochat\textsuperscript{\getAff{Stanford Research Comuting Center, Stanford University, Stanford, CA, US}},
Stephen Eglen\textsuperscript{\getAff{Department of Applied Mathematics and Theoretical Physics, University of
Cambridge, Cambridge, Cambridgeshire, GB}},
Tim Head\textsuperscript{\getAff{Wild Tree Tech, Zurich, CH}},
Tony Hirst\textsuperscript{\getAff{Department of Computing and Communications, The Open University, GB}}\\
\bigskip
\bigskip
* Corresponding author: daniel.nuest@uni-muenster.de\\
\end{flushleft}
% Please keep the abstract below 300 words
\section*{Abstract}
Containers are greatly improving computational science by packaging
software and data dependencies. In a scholarly context, transparency and
support of reproducibility are the largest drivers for using these
containers. It follows that choices that are made with respect to
building containers can make or break a worklow's reproducibility. The
build for the container image is often created based on the instructions
in a plain-text file. For example, one such container technology,
Docker, provides instructions using a \texttt{Dockerfile}. By following
the rules in this article researchers writing a \texttt{Dockerfile} can
effectively build and distribute containers.

% Please keep the Author Summary between 150 and 200 words
% Use first person. PLOS ONE authors please skip this step.
% Author Summary not valid for PLOS ONE submissions.

\linenumbers

% Use "Eq" instead of "Equation" for equation citations.
\hypertarget{introduction}{%
\section*{Introduction}\label{introduction}}
\addcontentsline{toc}{section}{Introduction}

With access to version control systems and only collaboration platforms
based on these, such as \href{https://github.com}{GitHub} or
\href{https://gitlab.com}{GitLab}, it has become increasingly easy to
not only share algorithms, but also instructions for executing research
workflows, including building and testing of the used software. The
publication of these artefacts are a response to the increasing
complexity of computer-based research, which is too complicated for
``papers'' to fully communicate the knowledge {[}1{]} but where the
actual knowledge is the full software environment that produced a result
{[}2{]}. A research compendium is a compilation of data, code, and
documentation that accompanies, or is itself, a scholarly publication
for increased transparency and reproducibility
(cf.~\url{https://research-compendium.science}). Within such a
compendium, it is suggested to include instructions for building
containers to package software and data dependencies {[}3{]}. By way of
providing these instructions along with thorough documentation and
application of common practices for scientific computing {[}4,5{]}, it's
much more likely both the author and third parties are able to reproduce
an analysis workflow. These instructions are human and machine readable
and the containers built from them are portable encapsulated snapshots
of a specific environment. For example, containers have been created to
host scientific notebooks {[}6{]}, and share reproducible workflows
(cf.~Rule 10 of {[}7{]}).

While there are several tutorials for using containers for reproducible
research {[}8--12{]}, there is no extensive examination for how to write
the actual instructions to create the containers. Several platforms for
facilitating reproducible research are built on top of containers
{[}13--17{]}, but they hide any complexity from the researcher. While
not everybody needs to understand and write Dockerfiles, as \emph{``the
number of unique research environments approximates the number of
researchers''} {[}17{]}, the researchers who do need to craft their own
environment definition should follow good practices. Such practices are
not part of generic Docker tutorials and are not present in published
Dockerfiles often used as templates. The differences between a helpful
stable Dockerfile and one that is misleading, prone to failure, and
contains potential obstacles, are not obvious, especially for
researchers who don't have software development experience.

While there are many different container technologies, this article
focuses on Docker {[}18{]}. The goal at hand is to write a
\texttt{Dockerfile} so that it best facilitates interactive development
and research as well as the higher goals of reproducibility and
preservation. A daily commitment to these general practices can ensure
that workflows are reproducible, and generation of a container is not an
afterthought triggered by a scientific publication (cf.~thoughts on
openness as an afterthought by {[}19{]} and on computational
reproducibility by {[}2{]}). By following the \emph{conventions} layed
out in these ten rules, authors ensure readability by others and ideally
subsequent reuse and collaboration.

\hypertarget{docker}{%
\subsection*{Docker}\label{docker}}
\addcontentsline{toc}{subsection}{Docker}

Docker is a container technology that is widely adopted and supported on
many platforms, and has become highly useful in science. They are
distinct from virtual machines (VM) or hypervisors as they do not
emulate hardware, and thus do not require the same system resources. To
build Docker containers, we write text files that follow a particular
format called \texttt{Dockerfile}s {[}20{]}. \texttt{Dockerfile}s are
machine- and human-readable recipes for creating containers. While
Docker was the original technology to support this format, other
container technologies have developed around the format (and thus
support it) including:
\href{https://podman.io/}{podman}/\href{https://github.com/containers/buildah}{buildah}
supported by RedHat,
\href{https://github.com/GoogleContainerTools/kaniko}{kaniko},
\href{https://github.com/genuinetools/img}{img}, and
\href{https://github.com/moby/buildkit}{buildkit}. The Singularity
container software {[}21{]} is optimized for high performance computing
and although it uses its own format, the \emph{Singularity recipe}, it
can import Docker containers directly from a Docker registry. Although
the Singularity recipe format is different, the rules here are
transferable to some extent. While some may argue for reasons to not
publish reproducibly (lack of time \& incentives and unfittingness for a
researcher's workflow, cf.~{[}22{]}) and there are substantial technical
challenges to maintain software and documentation, providing a
\texttt{Dockerfile} or pre-built Docker container should become an
increasingly easier task for the average researcher. If a researcher is
able to find or create containers or \texttt{Dockerfile}s to address
their most common use cases, then arguably it will not be extra work
after this initial set up (cf.~README of {[}23{]}). In fact, the
\texttt{Dockerfile} itself can be used as documentation to clearly show
from where data and code was derived, i.e.~downloaded or installed, and
consequently where a third party might obtain them again. Importantly,
the researcher should clearly state the platforms (by way of one or more
Docker containers) have been tested on.

\hypertarget{consider-tools-to-assist-with-dockerfile-generation}{%
\section*{1. Consider tools to assist with Dockerfile
generation}\label{consider-tools-to-assist-with-dockerfile-generation}}
\addcontentsline{toc}{section}{1. Consider tools to assist with
Dockerfile generation}

Writing a Dockerfile from scratch is not that simple, and even
``experts'' sometimes take shortcuts. Thus, it's a good strategy to
first look to tools that can help to generate a Dockerfile for you. Such
tools have likely thought about and implemented good practices, and they
may have added newer practices when reapplied at a later point in time.
It's always a good practice to start with your specific use case. You
would first want to determine if there is an already existing container
that you can use, and in this case, use it and add to your workflow
documentation instructions for doing so. As an example, you might be
doing some kind of interactive development. For interactive development
environments such as notebooks and development servers or databases, you
can readily find containers that come installed with the software that
you need. In the case that there isn't an existing container for your
needs, you might next look to well-maintained tools to help with
\texttt{Dockerfile} generation. Well maintained means that recent
container bases are used, likely from the official Docker library
{[}24{]}, to ensure that the container has the most recent security
fixes for the operating system in question. As an example, repo2docker
{[}16{]} is a tool that is maintained by Jupyter Labs that can help to
transform a repository in the format of some known kind of package
(notebook, Python, R, etc.) into a container. Such a package contains
well-defined files for defining software dependencies and versions,
which repo2docker can understand. As an example, we might install
\texttt{jupyter-repo2docker} and then run it against a repository with a
\texttt{requirements.txt} file, an indication of being a Python package
with the following command.

\begin{Shaded}
\begin{Highlighting}[]
\ExtensionTok{jupyter-repo2docker}\NormalTok{ https://github.com/norvig/pytudes}
\end{Highlighting}
\end{Shaded}

The resulting container image would install the dependencies listed in
the requirements file, along with providing an entrypoint to run a
notebook server to easily interact with any existing workflows in the
repository. A precaution that needs to be taken is that the default
command above will create a home for the current user, meaning that the
container itself wouldn't be ideal to share, but rather any researchers
interested in interaction with the code inside should build their own
container. For this reason, it's good practice to look at any help
command provided by the tool and check for configuration options for
user names, user ids, and similar. It's also recommended (if you are
able) to add custom labels to your container build to define metadata
for your analyses, if relevant. The container should be built to be
optimized for its use case, whether it is intended to be shared or used
by a single user.

Further assisting tools are \texttt{containerit} {[}25{]} and
\texttt{dockta} {[}26{]}. \texttt{containerit} automates the generation
of standalone Dockerfiles for workflows in R. It can provide a starting
point for users unfamiliar with writing Dockerfiles, or together with
other R packages provide a full image creation and execution process
without having to leave an R session. \texttt{dockta} supports multiple
programming languages and configurations files, just as
\texttt{repo2docker}, but attempts to create readable Dockerfiles
compatible with plain Docker and to improve user experience by cleverly
adjusting instructions to reduce build time.

However, in the case that a tool or interactive container environment is
not available, or their capabilities to configure specific resources or
install bespoke tools, you will likely need to write a
\texttt{Dockerfile} from scratch. In this case, proceed with following
the remaining rules to write your \texttt{Dockerfile}.

\hypertarget{use-versioned-and-automatically-built-base-images}{%
\section*{2. Use versioned and automatically built base
images}\label{use-versioned-and-automatically-built-base-images}}
\addcontentsline{toc}{section}{2. Use versioned and automatically built
base images}

It's good practice to use base images that are maintained by the Docker
library. While some organizations can be trustworthy in terms of
content, and reliable to update containers with security fixes (e.g.,
\texttt{rocker/r-ver}), for most individual accounts that build
containers, it's likely that containers will not be updated regularly.
It is even possible that images or Dockerfiles disappear. Therefore a
good understanding of how base images and image tags work is crucial, as
the tag that you choose has implications for your container build. You
should know how to use \texttt{FROM} statements to trace all the way up
to the original container base, an abstract base called
\texttt{scratch}. Doing this kind of trace is essential to be able to be
aware of all of the steps that were taken to generate your container and
who took them. If you want to build on a third party image, carefully
consider the author/maintainer and \emph{never} use an image without a
published Dockerfile. If you want to use an innofficial image, do safe a
copy of the Dockerfiles created by others. Alternatively copy relevant
instructions into your own Dockerfile, acknowledge the source in a
comment, and configure an automated build for your own image.

A tag like \texttt{latest} is good in that security fixes and updates
are likely present, however it is bad in that it's a moving target so
that it is more likely that an updated base could break your workflow.
Other tags that should be avoided are \texttt{dev}, \texttt{devel}, or
\texttt{nightly} that may provide possibly unstable versions of the
software. As an example, the Python container on Docker Hub will be
released with new versions of Python. If you build a container with
\texttt{python:latest} that only supports Python version 3, when the
base is updated to Python 4 it's likely that your software won't work as
expected. In this case, a tag like \texttt{python:3.5} might be a good
choice, where the third component of the version, i.e., \texttt{3.5.x},
ensures that security patches and bugfixes that won't break your code,
i.e., changes that are backwards compatible, are applied (cf.~semantic
versioninig, {[}27{]}). When you choose a base image, choose one with a
Linux distribution that supports the software stack you are using, and
also take into account the bases that are widely used by your community.
As an example, Ubuntu is heavily used for geospatial research, and so
the \texttt{rocker/geospatial} image would be a good choice for spatial
data science with R, or \texttt{jupyter/tensorflow-notebook} could be a
good choice for machine learning with Python. Especially when data is
involved, containers can grow quickly in size, and you should keep this
in mind. If you need to built smaller containers, consider a
\texttt{busybox} base or a multi-stage build, which allows to
selectively keep files from one build step to another. Smaller images
often are indicated by having \texttt{slim} or \texttt{minimal} as part
of the tag. Also take into account the libraries that you actually need.
Base images that have complex software installed (e.g.~machine learning
libraries, specific BLAS library) are helpful and fine to use. To not
rely on a third party or other individual to maintain the recipe and
container, you should copy the Dockerfile to your own repository, and
also provide a deployment for it via your own automated build. Trusting
that the third party \texttt{Dockerfile} will persist for your analyses
is a risky thing to do. For the new user, a potential list of existing
communities with regular builds and updates includes:

\begin{itemize}
\tightlist
\item
  \href{https://www.rocker-project.org/}{Rocker} for R {[}3{]}
\item
  \href{https://bioconductor.org/help/docker/}{Docker containers for
  Bioconductor} for bioinformatics
\item
  \href{https://hub.docker.com/_/neurodebian}{NeuroDebian images} for
  neuroscience {[}28{]}
\item
  {[}Jupyter Docker
  Stacks{]}(https://jupyter-docker-stacks.readthedocs.io/en/latest/index.html
  for Notebook-based computing
\item
  \href{https://hub.docker.com/r/taverna/taverna-server}{Taverna Server}
  for running Taverna workflows
\end{itemize}

For example, here is how we would use a base image \texttt{r-ver} with
tag \texttt{3.5.2} from the \texttt{rocker} organization on Docker Hub
(\texttt{docker.io})

\begin{verbatim}
FROM rocker/r-ver:3.5.2
\end{verbatim}

\hypertarget{use-formatting-and-favor-clarity}{%
\section{3. Use formatting and favor
clarity}\label{use-formatting-and-favor-clarity}}

It's good practice to think of the \texttt{Dockerfile} as being human
\emph{and} machine readable. This means that you should use indentation,
newlines, and comments to make your \texttt{Dockerfile}s well documented
and give them readability. Specifically, carefully indent commands and
their arguments to make clear what belongs together, especially when
connecting multiple commands in a \texttt{RUN} instruction with
\texttt{\&\&}. Use \texttt{\textbackslash{}} to break a single command
into multiple lines to ensure that no single line gets too long, and use
long versions of parameters for readability (e.g., \texttt{-\/-input}
instead of \texttt{-i}). When you need to change a directory, use
\texttt{WORKDIR}, because it not only creates the directory if it
doesn't exist but also persist the change across multiple \texttt{RUN}
instructions.

As you are writing the \texttt{Dockerfile}, put yourself in the shoes of
the next user of it. Are the choices made and commands being executed
clear, or is further comment warranted? When appropriate, you can add
links to discussion in version control, code repository issues, or
comments about specific usage and decisions. Dependencies can be grouped
in this fashion, which also makes it easier to spot changes if your
Dockerfile is under version control. It can even be helpful to include
comments about possibly more intuitive commands that did not work so you
do not fall into the same trap twice. If you find that you need to
remember an undocumented step, that's an indication that it should be
documented.

Labels are useful to provide a more structured form of documentation
about software. For example, software versions, maintainer contact
information, along with vendor specific metadata are commonly seen and
the OCI Image Format Specification provides some common label keys
{[}29{]}.

Comments can also include usage guides, which can be helpful to others
and a future you. As an example, it's helpful to provide commented lines
with \texttt{docker\ build} and \texttt{docker\ run} within the
\texttt{Dockerfile} to show how to build and run the image. These
comments can be especially relevant if volume mounts or ports are
important for using the container, and by putting them at the end of the
file, they are more likely to be seen, easy to copy paste after a
container build, and to be in consistent state compared to documentation
in another file. Here is an example of a commented section to show build
and usage.

\begin{verbatim}
# Build the images with
##> docker build --tag great_workflow .
# Run the image:
##> docker run --it --port 80:80 --volume ./input:/input --name gwf great_workflow
# Extract the data:
##> docker cp gwf:/output/ ./output
\end{verbatim}

If you were to discover a previously written \texttt{Dockerfile} and not
remember the container identifier you used, it would be represented in
the \texttt{Dockerfile} in the \texttt{-\/-name} CLI parameter.
Following a common coding aphorism, we might say \emph{``A Dockerfile
written three months ago may just as well have been written by someone
else''}. Here is an example of various kinds of comments that might be
useful for a \texttt{Dockerfile}:

\begin{verbatim}
# apt-get install specific version, use 'apt-cache madison <pkg>' to see available versions
RUN apt-get install python3-pandas=0.23.3+dfsg-4ubuntu1

# RUN command spreading several lines
RUN R -e 'getOption("repos")' && \
  install2.r \
    fortunes \
    here

# this library must be installed from source to get version newer than in sources

# following commands from instructions at LINK HERE
\end{verbatim}

Clarity is always more important than brevity. For example, if your
container uses a script to run a complex install routine, instead of
removing it from the container upon completion, which is commonly seen
in production Dockerfiles aiming at small image size, you should keep
the script in the container for a future user to inspect. Generally, try
to design the \texttt{RUN} statements so that each performs one scoped
action (e.g., download, compile, and install one tool).

If a \texttt{RUN} statement is longer than a page and requires
scrolling, this may be challenging for the next reader to digest, and
you should consider splitting it up into several. When you install
several system libraries, it's good practice to add comments about why
the dependencies are needed. This way, if a piece of software is removed
from the container, it will be easier to remove the system dependencies
that are no longer needed. If you intend to build the image more than
once (perhaps during development) and you can take advantage of build
caching to avoid execution of time-consuming instructions, e.g., install
from a remote resource or a file that gets cached. You should list
commands \emph{in order} of least likely to change to most likely to
change and use the \texttt{-\/-no-cache} flag to force a re-build of all
layers. A recommended ordering based on this metric might be:

\begin{enumerate}
\def\labelenumi{\arabic{enumi}.}
\tightlist
\item
  system libraries
\item
  language-specific libraries or modules 1. from repositories (binaries)
  2. from source
\item
  own software/scripts (if not mounted)
\item
  labels
\item
  \texttt{RUN}/\texttt{ENTRYPOINT}
\end{enumerate}

Finally, as a supplement to content inside the \texttt{Dockerfile}, it's
good practice to also write a section in a README alongside the
Dockerfile for exactly how to build, run, and otherwise interact with
the container. If a pre-build image is provided on Docker Hub, you
should direct the user to it.

\begin{itemize}
\tightlist
\item
  modularisation is a two-edged sword (point out
  \href{https://www.mankier.com/1/podman-build}{\texttt{podman}'s
  \texttt{\#include} directive}?) 
\item
  use a linter to follow common practices and consistency 
\end{itemize}

\hypertarget{pin-versions}{%
\section*{4. Pin versions}\label{pin-versions}}
\addcontentsline{toc}{section}{4. Pin versions}

The reproducibility of your \texttt{Dockerfile} is heavily dependent on
how well you define versions for dependencies inside.

\textbf{system libraries}

System library versions can largely come from the base image tag that
you choose to use (e.g., \texttt{ubuntu:18.04}) however, you can also
install specific versions of system packages with the respective package
manager. This practice is called version pinning (e.g., on apt on apt:
https://blog.backslasher.net/my-pinning-guidelines.html). Version
pinning is good practice if the version is relevant. For example, you
might want to demonstrate a bug, prevent a bug in an updated version, or
pin a working version if you suspect an update could lead to a problem.
Generally, system libraries are more stable than software modules
supporting analysis scripts, but in some cases they can be highly
relevant to your workflow. \emph{Installing from source} is a useful way
to install very specific versions, however it comes at the cost of
needing build libraries. Here is how you might go about listing the
currently installed versions of software on your system:

\begin{itemize}
\tightlist
\item
  Debian/Ubuntu: \texttt{dpkg\ -\/-list}
\item
  Alpine: \texttt{apk\ -vv\ info\textbar{}sort}
\item
  CentOS: \texttt{yum\ list\ installed} or \texttt{rpm\ -qa}
\end{itemize}

\textbf{extension packages and programming language modules}

In the case of needing to install packages or dependencies for a
specific language, package managers are a good option. Package managers
generally provide reliable mirrors or endpoints to download software,
and many packages are tested before release. Most package managers have
a client (CLI) interface that can easily be used from \texttt{RUN}
commands in your \texttt{Dockerfile}, along with various flavors of
``freeze'' commands that can output a text file listing all software
packages and versions
(cf.~https://markwoodbridge.com/2017/03/05/jupyter-reproducible-science.html
cited by {[}6{]}) The biggest risk with using package managers with
respect to \texttt{Dockerfile}s is outsourcing configuration to file
formats that are supported. As an example, here are configuration files
supported by commonly used languages in scientific programming:

\begin{itemize}
\tightlist
\item
  Python: \texttt{requirements\ .txt}, \texttt{xxx.yml} (Conda)
\item
  R: \texttt{DESCRIPTION}
\item
  Java: \texttt{mvn.xml}
\item
  JavaScript: \texttt{package.json} of \texttt{npm}
\end{itemize}

In all of the above, the user would be required to inspect the file or
the build to see what is installed. In some cases (e.g., conda) the
package manager is also able to make decisions about versions to
install, which is likely to lead to a non-reproducible build. For this
reason, in the case of having few packages, it is suggested to write the
install steps and versions directly into the \texttt{Dockerfile}. For
example, the \texttt{RUN} statement here:

\begin{Shaded}
\begin{Highlighting}[]
\ExtensionTok{RUN}\NormalTok{ pip install geopy==1.20.0 }\KeywordTok{&&} \KeywordTok{\textbackslash{}}
    \ExtensionTok{pip}\NormalTok{ install uszipcode==0.2.2}
\end{Highlighting}
\end{Shaded}

serves as more clear documentation in a \texttt{Dockerfile} than a
requirements.txt file that lists the same:

\begin{Shaded}
\begin{Highlighting}[]
\ExtensionTok{RUN}\NormalTok{ pip install -r requirements.txt}
\end{Highlighting}
\end{Shaded}

\begin{itemize}
\tightlist
\item
  how to do in Python (\texttt{==\ x.y.z})
\item
  do it in R with \texttt{versions} package, or by using MRAN (e.g.~via
  \texttt{r-ver} image)- JavaScript?
\item
  Julia: \texttt{add\ Package@1.0} \textgreater{}
  https://julialang.github.io/Pkg.jl/v1/managing-packages/\#Adding-packages-1
\item
  Use common command-line ready installation commands of programming
  languages

  \begin{itemize}
  \tightlist
  \item
    better readbiliy, potentially even performance
    (\texttt{RUN\ install2.r\ sp} instead of
    \texttt{RUN\ R\ -e\ "install.packages(\textquotesingle{}sf\textquotesingle{})"},
    although the latter is ``base R'')
  \end{itemize}
\end{itemize}

\hypertarget{mount-data-and-control-code}{%
\section*{5. Mount data and control
code}\label{mount-data-and-control-code}}
\addcontentsline{toc}{section}{5. Mount data and control code}

The addition of data or code into the container using
\texttt{ADD}/\texttt{COPY} is highly dependent on your use case. For
example, small data files and software that is essential for the
function of a container to distribute a reproducible analysis is
essential to add. However, if the container is intended to be used as an
interactive environment, you should refain from adding it to the
container, and instead bind the files from the outside of the container.
This also ensures that data persists when the container instance or
image is removed from your system. When you do this, in the case that
you will likely save as a root user from inside the container, you
should be careful about file permissions. To avoid problems when
mounting, Docker has provided a \texttt{-\/-user} option that you should
use. Without this option, there can be unexpected issues with
permissions. For example, writing a new file from inside the container
will be owned by user ``root'' on your host. In that many applications
generate log and other output files during runtime, be cautious about
what volumes are bound to the host. The other logical case to bind mount
code is when your data is very large (usually over hundreds of GB) or
has limited access. For example, a private dataset with personal health
information (PHI) should never be added to a container. In these cases,
you should provide clear instructions to the user in the README for how
to obtain actual or dummy data.

\begin{itemize}
\tightlist
\item
  easier access, does not require Docker knowledge by third parties to
  reause code and data
\item
  if you have a ``stable'' published software library, install it from
  source from the source code repo or from the software repository (so
  that users find the project in the future)
\item
  how to mount the data can be put in the example commands (see 4.)
\end{itemize}

\begin{verbatim}
docker run --volume ...
\end{verbatim}

\hypertarget{capture-environment-metadata}{%
\section*{6. Capture environment
metadata}\label{capture-environment-metadata}}
\addcontentsline{toc}{section}{6. Capture environment metadata}

Labels and build arguments can be very helpful to both provide metadata
and allow for customization of a build.

\begin{itemize}
\tightlist
\item
  \textbf{labels}
\end{itemize}

Labels serve as structured metadata that can be exposed by APIs (e.g.,
https://microbadger.com/labels) along with tools to inspect the
container binaries (e.g., \texttt{docker\ inspect}). Namespaced
conventions have been developed (http://label-schema.org/rc1/
respectively https://github.com/opencontainers/image-spec) to help
standardize across container tools. Important metadata attributes to
include as labels would include any of the following:

\begin{itemize}
\tightlist
\item
  Research organizations (https://ror.org/)
\item
  Funding agency/grant number
\item
  DOI of research compendium
\item
  Author DOIs (e.g., Zenodo or other)
\item
  repository link where Dockerfile is published
\item
  author (\texttt{MAINTAINER} is deprecated) and contact (e.g.~email,
  project website)
\item
  license
\end{itemize}

There are several projects (https://codemeta.github.io/,
https://citation-file-format.github.io/) focused around metadata
citation, and you should consider using one if it seems appropriate for
your use case.

\begin{itemize}
\tightlist
\item
  \textbf{build arguments}
\end{itemize}

Build arguments can provide metadata and also allow for customization of
a build. As an example, the following build argument would default to
\texttt{1.0.0} but allow the user to change it with
\texttt{-\/-build-arg\ MYVERSION=2.0.0}

\begin{verbatim}
ARG MYVERSION=1.0.0
\end{verbatim}

Along with specifying versions, environment, or timestamp, build
arguments can be useful to provide the context of the build (e.g.,
production, development, automated or not). Examples of build arguments
that are useful to include to describe a container are:

\begin{itemize}
\tightlist
\item
  git commit hash
\item
  date and timestamp
\end{itemize}

\hypertarget{enable-interactive-usage-and-one-click-execution}{%
\section*{7. Enable interactive usage and one-click
execution}\label{enable-interactive-usage-and-one-click-execution}}
\addcontentsline{toc}{section}{7. Enable interactive usage and one-click
execution}

The success of a container can sometimes come down to how well it
exposes it's usage and software to the user: a function that comes down
to the \texttt{CMD} and \texttt{ENTRYPOINT}. By using \texttt{CMD} and
\texttt{ENTRYPOINT}, you can use automated and manual testing to ensure
that it is possible to run the container interactively \emph{and} given
that execution is done incorrectly or has an error, a suitable help or
error message is shown to assist the user. A possible weakness with
using containers is the limitation on only providing one entrypoint,
however tools (e.g., The Scientific Filesystem REF
academic.oup.com/gigascience/article/7/5/giy023/4931737) have been
developed to expose multiple entrypoints, environments, help messages,
labels, and even install sequences. It's considered good practice to
have an entrypoint that meets the user expectations. For example, a
container known to be a workflow should execute the workflow, or spit
out usage for how to do so. An interactive container should spin up an
analysis environment, ensuring to print to the screen the ports that are
being used (that need to be bound to the host) along with any login
credentials needed. A workflow that does not support headless execution
is arguably not ideal for a container, and you should consider switching
tools. Using a browser to expose a user interface (e.g, RStudio,
Jupyter, web tools) on a particular port is a commonly done practice,
and for systems without native \texttt{x11} support, \texttt{x11docker}
is recommended {[}REF{]}. Otherwise, you would need to clearly document
how the user can start and control the graphical interfaces provided
inside. However, headless execution is ideal in that it can be tested
headlessly, either by starting or interacting with the container and
checking for successful responses (e.g., 200) from endpoints provided,
or by using a controller such as Selenium {[}REF{]}. To support both
``one click execution'' and to allow for custom configuration, it's
helpful to provide users with a configuration file, or allow for
settings to be expressed via environment variables {[}30{]}, or special
Docker-based wrappers such as Kliko {[}31{]}

\begin{itemize}
\tightlist
\item
  may also use the same \texttt{Dockerfile} for different purposes,
  e.g.~include an app (e.g.~Shiny) for interactive use by user
\item
  see also Rule 7: Build a pipeline in {[}6{]}
\item
  \emph{what user should run within the Dockerfile?} 
\end{itemize}

\hypertarget{establish-templates-for-new-projects}{%
\section{8. Establish templates for new
projects}\label{establish-templates-for-new-projects}}

It's likely going to be the case that you develop interactive
environments, servers, or workflows that are similar in nature. In this
case, you should consider adopting a standard workflow that will give
you a clean slate for a new project. As an example, cookie cutter
templates (e.g., https://github.com/cookiecutter/cookiecutter) and
project starter kits can provide files, directory organization, and
build instructions that include basic steps for getting started. A good
project template should get you started with a template for
documentation, setting up testing via continuous integration, building a
container, and even choosing a license. In the case of using a common
\texttt{Dockerfile} or base, Docker's build caching will take shared
lines into account and speed up build time. It also must be noted that
developing or working on projects with containers, period, allows you to
easily switch between isolated project environments. At most, if a port
is shared by two projects, you would simply need to stop the container
and restart when you are ready to work again. In the case of more
complex web applications that require application and web servers,
databases, and workers or messenging, the entire infrastructure can
easily be brought up or down with a tool like \texttt{docker-compose}
{[}REF{]}. \texttt{docker-compose} also allows definition of services
via it's own \texttt{docker-compose.yml}, ranging from mounted volumes,
to permissions, environment, and ports. Part of this project template
should be a protocol for publishing the container image to a suitable
container registry, and taking into consideration of how the code can be
given a DOI or proper publication (e.g., Zenodo, Journal of Open Source
Software).

\hypertarget{publish-a-dockerfile-per-project-in-a-code-repository-with-version-control}{%
\section*{9. Publish a Dockerfile per project in a code repository with
version
control}\label{publish-a-dockerfile-per-project-in-a-code-repository-with-version-control}}
\addcontentsline{toc}{section}{9. Publish a Dockerfile per project in a
code repository with version control}

The \texttt{Dockerfile} is a plain text-based format and thus fits
nicely into a version control system. You should tend to include
\texttt{Dockerfile}s alongside your code as a way to quickly build your
software, and show visitors of the repository how it's built (qoute
Abelson? see also
https://www.quora.com/How-true-is-Programs-are-meant-to-be-read-by-humans-and-only-incidentally-for-computers-to-execute)/,
and as mentioned previously, the \texttt{Dockerfile} itself can then
link back to the repository. Take advantage of git commit messages to
keep a record of notes alongside changes to your files. Versioning and
publishing on a collaboration platform is also beneficial in that is
enables collaboration and makes it easy to solicit feedback. Version
control providers (e.g., GitHub, GitLab) also make it easy to connect to
or use embedded continuous integration (CI) services to test and deploy
your containers and software. While it's not always possible to keep
datasets alongside code, in the case of smaller ones, it can be very
useful to keep your \texttt{Dockerfile} alongside both the code and data
(cf.~research compendium concept?). While there are exceptions to the
rule (cf.~{[}32{]}), it's generally a simple and clear approach to
provide one Dockerfile per project. If you find that you need to provide
more than one, consider if it's possible to use build arguments to flip
between states (e.g., development vs.~production) or if it's possible to
separate tools into different repositories. In fact, using and
publishing a \texttt{Dockerfile} to create a container has been shown to
increase chances of preservation of your work (cf.~{[}33{]})

Consult Ten Simple Rules paper by Perez-Riverol et al.~{[}34{]}

\hypertarget{use-the-container-daily-rebuild-the-image-weekly}{%
\section*{10. Use the container daily, rebuild the image
weekly}\label{use-the-container-daily-rebuild-the-image-weekly}}
\addcontentsline{toc}{section}{10. Use the container daily, rebuild the
image weekly}

A container that is not regularly rebuilt risks not getting essential
security updates, and serving as a risk for its users. If you are using
a container built by a \texttt{Dockerfile} regularity in your work, it's
recommended to set up continuous deployment, or a strategy to ensure
that the container is regularly built and tested. While there is no
hardened rule, it would be good practice to rebuild a container every
once or two weeks that you use on a daily basis. In an idealistic world,
containers that we built years ago should rebuild seamlessly, but this
isn't the case, especially with rapidly changing technology relevant to
machine learning and data science. It can almost be guaranteed that the
longer that you wait to recompile the image, the harder it will become.
If you are using an interactive container and find that you need to
manually installed a package or change a parameter, it would be best
practice to add this dependency to the container and rebuild it. If the
container is linked to an automated build provided by a version control
service, simply pushing the \texttt{Dockerfile} to the service can
easily trigger the build. Docker images and layers can exhaust system
resources, so it's good practice to regularly do
\texttt{system\ prune\ -\/-all} that will clean up stopped and unused
containers, volumes, and networks. If you want to regularly prune unused
container resources, you might even consider adding a cron job {[}REF{]}
to do it. After a prune is performed, it would follow naturally to
rebuild a container for local usage, or to pull it again from a newly
built registry image. In the case of needing more setup or configuration
than building the container, it's good practice to provide a
\texttt{Makefile} alongside your container. You generally should not
rely on undocumented steps or manual, extra commands to be run on the
local machine. A fully scripted configuration makes it easier for both
you and future users, and can increase trust in your workflow. Finally,
for a sanity check it can be helpful to ask a colleague to interact with
your workflow or container.

in your regular work, it is the only way to make sure it is really
stable (cf.~Marwick's ``this container is the only way I have ever run
this workflow'') - no showstopper for using UIs (web-based,
e.g.~Jupyter, RStudio, but also \texttt{x11docker}) - \textbf{Don't
replicate environment configuration outside of the Dockerfile for
convenience} - avoid an untidy laboratory in practice behind a shiny
appearance of a \texttt{Dockerfile} - you can install interactive UIs as
part of the Dockerfile and use them just like Desktop UIs (Jupyter,
RStudio, use )

\hypertarget{example-dockerfiles}{%
\section{Example Dockerfiles}\label{example-dockerfiles}}

To demonstrate the 10 rules, we have a git repository with example
\texttt{Dockerfile}s, some of which we took from public repositories and
updated to adhere to the rules (\texttt{Dockerfile.before} and
\texttt{Dockerfile.after}).

\hypertarget{conclusion}{%
\section*{Conclusion}\label{conclusion}}
\addcontentsline{toc}{section}{Conclusion}

Reproducibility in research is an endeavour of best efforts, not about
achieving the perfect solution, as that is probably not achievable or
changing over time. This article provides guidance for using
\texttt{Dockerfile}s in computational/computer-based research to work
towards a ``time capsule'' (see
https://twitter.com/DougBlank/status/1135904909663068165?s=09) which,
given some expertise and the right tools, can be used to come as close
as possible to the original workflow. Even if such a capsule decays over
time, the effort to create it provides incredibly useful and valuable
transparency for the project. The effort should also be valued by others
and may change the way scholars collaborate and communicate (cf.~notion
of ``preproducibility'' by {[}35{]}) So please, don't go insane with
writing \texttt{Dockerfile}s, but be realistic about what might break
and what is unlikely to break, and make a best effort. In a similar
vein, these rules may be broken if another way works better for
\emph{your use case}. The rules in this article help you mastering the
\texttt{Dockerfile} format for research and provide a solid basis for
engaging in more complex but also in simpler assisted usage of
containers (cf.~Box: Aumatic Generation). Corner cases aside, share and
exchange your \texttt{Dockerfile} freely and collaborate in your
community to spread the knowledge about containers as a tool for
research. Together you can develop common practices or even shared base
images (exemplified by communities listed in Rule 1).

\hypertarget{acknowledgements}{%
\section*{Acknowledgements}\label{acknowledgements}}
\addcontentsline{toc}{section}{Acknowledgements}

o2r by DFG

\hypertarget{contributions}{%
\section*{Contributions}\label{contributions}}
\addcontentsline{toc}{section}{Contributions}

DN conceived the idea, wrote the first darft, contributed to all rules.
VS converted first draft bullet points into full sentences, contributed
to all rules. SE contributed to the overall structure and selected
rules. TH contributed to the rule structure and particularly rule 1. THi
gave extensive feedback on an early draft.

\hypertarget{references}{%
\section*{References}\label{references}}
\addcontentsline{toc}{section}{References}

\hypertarget{refs}{}
\leavevmode\hypertarget{ref-marwick_how_2015}{}%
1. Marwick B. How computers broke science -- and what we can do to fix
it {[}Internet{]}. The Conversation. 2015. Available:
\url{https://theconversation.com/how-computers-broke-science-and-what-we-can-do-to-fix-it-49938}

\leavevmode\hypertarget{ref-donoho_invitation_2010}{}%
2. Donoho DL. An invitation to reproducible computational research.
Biostatistics. 2010;11: 385--388.
doi:\href{https://doi.org/10.1093/biostatistics/kxq028}{10.1093/biostatistics/kxq028}

\leavevmode\hypertarget{ref-boettiger_introduction_2017}{}%
3. Boettiger C, Eddelbuettel D. An Introduction to Rocker: Docker
Containers for R. The R Journal. 2017;9: 527--536.
doi:\href{https://doi.org/10.32614/RJ-2017-065}{10.32614/RJ-2017-065}

\leavevmode\hypertarget{ref-wilson_best_2014}{}%
4. Wilson G, Aruliah DA, Brown CT, Hong NPC, Davis M, Guy RT, et al.
Best Practices for Scientific Computing. PLOS Biology. 2014;12:
e1001745.
doi:\href{https://doi.org/10.1371/journal.pbio.1001745}{10.1371/journal.pbio.1001745}

\leavevmode\hypertarget{ref-wilson_good_2017}{}%
5. Wilson G, Bryan J, Cranston K, Kitzes J, Nederbragt L, Teal TK. Good
enough practices in scientific computing. PLOS Computational Biology.
2017;13: e1005510.
doi:\href{https://doi.org/10.1371/journal.pcbi.1005510}{10.1371/journal.pcbi.1005510}

\leavevmode\hypertarget{ref-rule_ten_2019}{}%
6. Rule A, Birmingham A, Zuniga C, Altintas I, Huang S-C, Knight R, et
al. Ten simple rules for writing and sharing computational analyses in
Jupyter Notebooks. PLOS Computational Biology. 2019;15: e1007007.
doi:\href{https://doi.org/10.1371/journal.pcbi.1007007}{10.1371/journal.pcbi.1007007}

\leavevmode\hypertarget{ref-sandve_ten_2013}{}%
7. Sandve GK, Nekrutenko A, Taylor J, Hovig E. Ten Simple Rules for
Reproducible Computational Research. PLoS Comput Biol. 2013;9: e1003285.
doi:\href{https://doi.org/10.1371/journal.pcbi.1003285}{10.1371/journal.pcbi.1003285}

\leavevmode\hypertarget{ref-nust_author_2017}{}%
8. Nüst D. Author Carpentry : Docker for reproducible research
{[}Internet{]}. Author Carpentry : Docker for reproducible research.
2017. Available:
\url{https://nuest.github.io/docker-reproducible-research/}

\leavevmode\hypertarget{ref-chapman_reproducible_2018}{}%
9. Chapman P. Reproducible data science environments with Docker Phil
Chapman's Blog {[}Internet{]}. 2018. Available:
\url{https://chapmandu2.github.io/post/2018/05/26/reproducible-data-science-environments-with-docker/}

\leavevmode\hypertarget{ref-ropensci_labs_r_2015}{}%
10. rOpenSci Labs. R Docker tutorial {[}Internet{]}. 2015. Available:
\url{https://ropenscilabs.github.io/r-docker-tutorial/}

\leavevmode\hypertarget{ref-udemy_docker_2019}{}%
11. Udemy, Zhbanko V. Docker Containers for Data Science and
Reproducible Research {[}Internet{]}. Udemy. 2019. Available:
\url{https://www.udemy.com/course/docker-containers-data-science-reproducible-research/}

\leavevmode\hypertarget{ref-psomopoulos_lesson_2017}{}%
12. Psomopoulos FE. Lesson "Docker and Reproducibility" in Workshop
"Reproducible analysis and Research Transparency" {[}Internet{]}.
Reproducible analysis and Research Transparency. 2017. Available:
\url{https://reproducible-analysis-workshop.readthedocs.io/en/latest/8.Intro-Docker.html}

\leavevmode\hypertarget{ref-brinckman_computing_2018}{}%
13. Brinckman A, Chard K, Gaffney N, Hategan M, Jones MB, Kowalik K, et
al. Computing environments for reproducibility: Capturing the ``Whole
Tale''. Future Generation Computer Systems. 2018;
doi:\href{https://doi.org/10.1016/j.future.2017.12.029}{10.1016/j.future.2017.12.029}

\leavevmode\hypertarget{ref-code_ocean_2019}{}%
14. Code Ocean {[}Internet{]}. 2019. Available:
\url{https://codeocean.com/}

\leavevmode\hypertarget{ref-simko_reana_2019}{}%
15. Šimko T, Heinrich L, Hirvonsalo H, Kousidis D, Rodríguez D. REANA: A
System for Reusable Research Data Analyses. EPJ Web of Conferences.
2019;214: 06034.
doi:\href{https://doi.org/10.1051/epjconf/201921406034}{10.1051/epjconf/201921406034}

\leavevmode\hypertarget{ref-jupyter_binder_2018}{}%
16. Jupyter P, Bussonnier M, Forde J, Freeman J, Granger B, Head T, et
al. Binder 2.0 - Reproducible, interactive, sharable environments for
science at scale. Proceedings of the 17th Python in Science Conference.
2018; 113--120.
doi:\href{https://doi.org/10.25080/Majora-4af1f417-011}{10.25080/Majora-4af1f417-011}

\leavevmode\hypertarget{ref-nust_opening_2017}{}%
17. Nüst D, Konkol M, Pebesma E, Kray C, Schutzeichel M, Przibytzin H,
et al. Opening the Publication Process with Executable Research
Compendia. D-Lib Magazine. 2017;23.
doi:\href{https://doi.org/10.1045/january2017-nuest}{10.1045/january2017-nuest}

\leavevmode\hypertarget{ref-wikipedia_contributors_docker_2019}{}%
18. Wikipedia contributors. Docker (software) {[}Internet{]}. Wikipedia.
2019. Available:
\url{https://en.wikipedia.org/w/index.php?title=Docker_(software)\&oldid=928441083}

\leavevmode\hypertarget{ref-chen_open_2019}{}%
19. Chen X, Dallmeier-Tiessen S, Dasler R, Feger S, Fokianos P, Gonzalez
JB, et al. Open is not enough. Nature Physics. 2019;15: 113.
doi:\href{https://doi.org/10.1038/s41567-018-0342-2}{10.1038/s41567-018-0342-2}

\leavevmode\hypertarget{ref-docker_inc_dockerfile_2019}{}%
20. Docker Inc. Dockerfile reference {[}Internet{]}. Docker
Documentation. 2019. Available:
\url{https://docs.docker.com/engine/reference/builder/}

\leavevmode\hypertarget{ref-kurtzer_singularity_2017}{}%
21. Kurtzer GM, Sochat V, Bauer MW. Singularity: Scientific containers
for mobility of compute. PLOS ONE. 2017;12: e0177459.
doi:\href{https://doi.org/10.1371/journal.pone.0177459}{10.1371/journal.pone.0177459}

\leavevmode\hypertarget{ref-boettiger_introduction_2015}{}%
22. Boettiger C. An Introduction to Docker for Reproducible Research.
SIGOPS Oper Syst Rev. 2015;49: 71--79.
doi:\href{https://doi.org/10.1145/2723872.2723882}{10.1145/2723872.2723882}

\leavevmode\hypertarget{ref-marwick_madjebebe_2015}{}%
23. Ben Marwick. 1989-excavation-report-Madjebebe. 2015;
doi:\href{https://doi.org/10.6084/m9.figshare.1297059}{10.6084/m9.figshare.1297059}

\leavevmode\hypertarget{ref-docker_inc_official_2019}{}%
24. Docker Inc. Official Images on Docker Hub {[}Internet{]}. Docker
Documentation. 2019. Available:
\url{https://docs.docker.com/docker-hub/official_images/}

\leavevmode\hypertarget{ref-nust_containerit_2019}{}%
25. Nüst D, Hinz M. Containerit: Generating Dockerfiles for reproducible
research with R. Journal of Open Source Software. 2019;4: 1603.
doi:\href{https://doi.org/10.21105/joss.01603}{10.21105/joss.01603}

\leavevmode\hypertarget{ref-stencila_dockta_2019}{}%
26. Stencila. Stencila/dockta {[}Internet{]}. Stencila; 2019. Available:
\url{https://github.com/stencila/dockta}

\leavevmode\hypertarget{ref-preston-werner_semantic_2013}{}%
27. Preston-Werner T. Semantic Versioning 2.0.0 {[}Internet{]}. Semantic
Versioning. 2013. Available: \url{https://semver.org/}

\leavevmode\hypertarget{ref-halchenko_open_2012}{}%
28. Halchenko YO, Hanke M. Open is Not Enough. Let's Take the Next Step:
An Integrated, Community-Driven Computing Platform for Neuroscience.
Frontiers in Neuroinformatics. 2012;6.
doi:\href{https://doi.org/10.3389/fninf.2012.00022}{10.3389/fninf.2012.00022}

\leavevmode\hypertarget{ref-opencontainers_image-spec_2017}{}%
29. Opencontainers. Opencontainers/image-spec v1.0.1 - Annotations
{[}Internet{]}. GitHub. 2017. Available:
\url{https://github.com/opencontainers/image-spec/blob/v1.0.1/annotations.md}

\leavevmode\hypertarget{ref-knoth_reproducibility_2017}{}%
30. Knoth C, Nüst D. Reproducibility and Practical Adoption of GEOBIA
with Open-Source Software in Docker Containers. Remote Sensing. 2017;9:
290. doi:\href{https://doi.org/10.3390/rs9030290}{10.3390/rs9030290}

\leavevmode\hypertarget{ref-molenaar_klikoscientific_2018}{}%
31. Molenaar G, Makhathini S, Girard JN, Smirnov O. Kliko---The
scientific compute container format. Astronomy and Computing. 2018;25:
1--9.
doi:\href{https://doi.org/10.1016/j.ascom.2018.08.003}{10.1016/j.ascom.2018.08.003}

\leavevmode\hypertarget{ref-kim_bio-docklets_2017}{}%
32. Kim B, Ali TA, Lijeron C, Afgan E, Krampis K. Bio-Docklets:
Virtualization Containers for Single-Step Execution of NGS Pipelines.
bioRxiv. 2017; 116962.
doi:\href{https://doi.org/10.1101/116962}{10.1101/116962}

\leavevmode\hypertarget{ref-emsley_framework_2018}{}%
33. Emsley I, De Roure D. A Framework for the Preservation of a Docker
Container International Journal of Digital Curation. International
Journal of Digital Curation. 2018;12.
doi:\href{https://doi.org/10.2218/ijdc.v12i2.509}{10.2218/ijdc.v12i2.509}

\leavevmode\hypertarget{ref-perez-riverol_ten_2016}{}%
34. Perez-Riverol Y, Gatto L, Wang R, Sachsenberg T, Uszkoreit J,
Leprevost F da V, et al. Ten Simple Rules for Taking Advantage of Git
and GitHub. PLOS Computational Biology. 2016;12: e1004947.
doi:\href{https://doi.org/10.1371/journal.pcbi.1004947}{10.1371/journal.pcbi.1004947}

\leavevmode\hypertarget{ref-stark_before_2018}{}%
35. Stark PB. Before reproducibility must come preproducibility
{[}Internet{]}. Nature. 2018.
doi:\href{https://doi.org/10.1038/d41586-018-05256-0}{10.1038/d41586-018-05256-0}

\nolinenumbers


\end{document}

